\section{Related Work}
This section reviews the literature most closely related to wearable navigability estimation using privacy-preserving thermal sensing and multi-view perception.
\subsection{Wearable Navigation and Traversability Estimation}
Existing wearable navigation systems estimate traversable space with sensors
that are accurate but hardware-heavy and privacy-invasive. Most rely on RGB
cameras~\cite{11540071}, depth sensors, or LiDAR~\cite{Liu_2021_ICCV} to
detect obstacles and identify safe directions for movement. RGB-based methods
infer navigability from visual
appearance~\cite{muDynamicObstacleAvoidance2024}, whereas depth- and
LiDAR-based systems rely on explicit geometric measurements of the
environment. These sensing requirements increase hardware complexity and raise
privacy concerns for continuous wearable
deployment~\cite{rinta-homiHowLowCan2021}.

\subsection{Privacy-Preserving and Ultra-Low-Resolution Thermal Sensing}

Privacy-preserving thermal sensing has emerged as a promising alternative to RGB-based perception for wearable navigation, enabling continuous environmental sensing while reducing concerns associated with capturing identifiable imagery~\cite{papagianopoulos2023obstacle}. Recent work has demonstrated the feasibility of thermal-based obstacle detection using standard-resolution thermal cameras. In particular, DareDevil~\cite{gaoDeployablePrivacyPreservingThermalBased2025} leverages a paired RGB camera and a pretrained YOLOv8 detector to automatically label thermal frames, eliminating the need for manual annotation while performing thermal-only obstacle detection at inference. However, DareDevil operates on imagery from a single standard-resolution thermal camera ($640\times480$) and formulates the task as binary obstacle detection. It therefore does not investigate directional navigability estimation, multi-camera thermal perception, or the challenges posed by ultra-low-resolution thermal sensing.

% DareDevil~\cite{gaoDeployablePrivacyPreservingThermalBased2025} is the closest
% precedent to our labeling pipeline, auto-labeling thermal frames from a paired
% RGB camera through a pretrained YOLOv8 detector and thereby avoiding manual
% annotation. However, DareDevil collapses all
% object classes into a single binary obstacle label and operates on colorized
% thermal frames from a single standard-resolution thermal camera at inference
% time. It does not address the extreme spatial compression of a
% $32\times24$-pixel array, multi-camera fusion, or robustness to sensor
% placement.

Ultra-low-resolution thermal arrays further improve privacy, cost, and power efficiency, but introduce substantially greater perception challenges due to their extremely limited spatial resolution. Existing work has therefore focused on tasks such as occupancy detection~\cite{rinta-homiHowLowCan2021}, human monitoring~\cite{rezaeiUnobtrusiveHumanActivity2023}, and motion estimation rather than scene understanding or navigability. For example, Mohsen et al.~\cite{mohsenDeepFusionUltraLowResolution2025} combine thermal imagery with gyroscope measurements to compensate for the limited visual information available for odometry.
% while Naser et al.~\cite{naser2022calibration} assume a fixed, calibrated sensor geometry. 
Whether ultra-low-resolution thermal sensing alone can support directional navigability estimation in a wearable multi-camera setting remains largely unexplored.

% tete line of work uses ultra-low-resolution thermal sensors (e.g.,
% $32\times24$ arrays, on the order of \$30 per unit) as a low-power,
% lighting-invariant alternative to RGB and depth sensing. Such sensors are too
% coarse to resolve identifiable facial features, but their reduced spatial
% detail also degrades the features on which downstream perception models rely.
% Prior work compensates for this by fusing thermal imagery with gyroscope
% readings for motion estimation~\cite{mohsenDeepFusionUltraLowResolution2025}.
% Naser~\textit{et al.}~\cite{naser2022calibration} separately showed that most
% low-resolution thermal pipelines still assume a fixed, calibrated sensor
% geometry. In contrast, this work targets discrete navigability classification
% rather than continuous odometry, and removes the reliance on precise geometric
% calibration at inference that these systems retain.

% Removed, not very relevant
% \subsection{Automatic Label Generation}
% Manually annotating walkable regions is expensive and does not scale to
% real-world indoor environments. Kim~\textit{et al.}~\cite{kim2024traversability}
% addressed this for egocentric video with an automated annotation strategy that
% derives traversability labels without manual pixel-level annotation. At larger
% scale, CityWalker~\cite{liuCityWalkerLearningEmbodied2025} trains an embodied
% urban navigation policy on over 2{,}000 hours of in-the-wild web video,
% replacing teleoperated demonstrations with pseudo-labels derived from visual
% odometry. We apply the same philosophy in a cross-modal setting, where
% RGB-derived object detections supervise a thermal-only model solely during
% data collection and the RGB sensor is entirely absent at deployment.

\subsection{Multi-View Sensor Fusion and Placement Robustness}

Multi-view perception systems combine observations from multiple cameras to increase environmental coverage and improve perception quality~\cite{zhuAutonomousRobotNavigation2020,olagokeLiteratureSurveyMultiCamera2020}. Existing systems typically aggregate information from multiple cameras within a unified perception pipeline~\cite{Zheng_2023_WACV}, although the choice of where and how to combine observations varies across applications. In contrast to prior work, this paper investigates how different fusion strategies perform for directional navigability estimation from wearable ultra-low-resolution thermal sensing, while also evaluating their robustness to changes in camera placement.

Wearable multi-camera platforms introduce an additional deployment challenge: the physical arrangement of sensors may change between training and deployment as cameras are repositioned or reattached. Existing multi-camera perception systems are typically designed for fixed camera configurations, an assumption that simplifies geometric reasoning and sensor fusion but is difficult to maintain in wearable platforms~\cite{kuoRedesigningSLAMArbitrary2020}. Consequently, changing the mapping between physical cameras and network input channels can degrade the performance of fusion models that exploit camera-specific spatial relationships. In this work, we evaluate the robustness of representative fusion strategies to camera permutations and investigate lightweight mechanisms for maintaining performance under changing sensor arrangements.

Taken together, the literature reviewed above has not established whether wearable ultra-low-resolution multi-view thermal sensing can reliably support directional navigability estimation under realistic deployment conditions.
% Multi-camera perception systems typically assume a fixed mapping between
% physical sensor position and input channel, an assumption that breaks when
% sensors shift or are reattached in a different order on a wearable rig.
% Pixel-level, decision-level, and feature-level fusion are well studied in the
% sensor fusion literature, but little work examines how these strategies behave
% when the camera-to-slot assignment is perturbed, particularly at the small
% camera counts and extreme resolutions considered here. This also relates to
% cross-resolution knowledge transfer, where a high-resolution signal supervises
% a model on coarser inputs~\cite{qiMultiScaleAlignedDistillation2021}. Our gap
% is larger still, since the supervisory RGB detector and the deployed
% $32\times24$ thermal arrays share no feature space and no geometric alignment
% at inference time. 

%We contribute a direct comparison of three fusion strategies for an ultra-low-resolution, three-camera array, together with an explicit treatment of camera-slot permutation robustness.
